% !TEX root = ../notes.tex

\documentclass[../notes.tex]{subfiles}

\begin{document}

\section{September 17}
Today, we discuss dualizable presentable $\infty$-categories.

\subsection{Tensor Products}
Let's return to tensor products. We are hunting for a symmetric monoidal structure on $\op{Pr}^L_\kappa$.
\begin{theorem}
	There is a unique functor
	\[-\otimes-\colon\mathrm{Pr}^L\times\mathrm{Pr}^L\to\mathrm{Pr}^L\]
	such that
	\[\op{Hom}_{\mathrm{Pr}^L}(\mc C\otimes\mc D,\mc E)=\op{Bilin}_{\mathrm{Pr}^L}(\mc C\times\mc D,\mc E).\]
	Here, $\op{Bilin}$ means that we have functors $\mc C\times\mc D\to\mc E$ which preserve colimits in each variable.
\end{theorem}
\begin{proof}
	We will directly construct $\otimes$. For any $\mc C\in\op{Pr}^L$, we may choose a presentation $\mc C=\langle S\rangle\left[R^{-1}\right]$, where $S$ is a small $\infty$-category, and $R$ is a collection of morphisms. (Indeed, such a presentation exists by writing out how $\mc C$ is compactly generated. For example, $\langle S\rangle$ is the ``free presentable'' category on $S$, which is $\op{PSh}(S;\mathrm{Ani})$.) Accordingly, given $\mc C$ and $\mc D$ with presentations $\mc C=\langle S_{\mc C}\rangle\left[R_{\mc C}^{-1}\right]$ and $\mc D=\langle S_{\mc D}\rangle\left[R_{\mc D}^{-1}\right]$, we define
	\[\mc C\otimes\mc D\coloneqq\langle S_{\mc C}\times S_{\mc D}\rangle\left[(R_{\mc C}\times1\cup1\times R_{\mc D})^{-1}\right].\]
	One can check the universal property from here fairly directly.
\end{proof}
\begin{remark}
	There is some intuition that colimits should give our additive structure, which motivates our definition of $\op{Bilin}$.
\end{remark}
\begin{remark}
	In fact, we have the adjunction property
	\[\op{Hom}_{\mathrm{Pr}^L}(\mc C\otimes\mc D,\mc E)=\op{Hom}_{\mathrm{Pr}^L}(\mc C,\op{Fun}(\mc D,\mc E)).\]
	Indeed, it is enough to check that the right-hand set is $\op{Bilin}_{\op{Pr}^L}(\mc C\times\mc D,\mc E)$, which we see by currying.
\end{remark}
\begin{remark}
	If $\mc C$ and $\mc D$ are $\kappa$-presentable, then $\mc C\otimes\mc D$ is also $\kappa$-present\-able. Indeed, being $\kappa$-present\-able allows  us to write $\mc C=\langle\mc C^\kappa\rangle\left[R_{\mc C}^{-1}\right]$ and $\mc D=\langle\mc D^\kappa\rangle\left[R_{\mc D}^{-1}\right]$, so our explicit construction of $\mc C\otimes\mc D$ tells us that it is generated by $\kappa$-compact objects, given by inverting only morphisms between $\kappa$-compact objects. In fact, one can check that $-\otimes-$ restricts to a functor $\mathrm{Pr}^L_\kappa\times\mathrm{Pr}^L_\kappa\to\mathrm{Pr}^L_\kappa$: one only needs to check that it sends compact objects to compact objects, which can be done via the presentation.
\end{remark}
\begin{example} \label{ex:dualizable-object}
	Suppose $\mc C=\langle S\rangle$ so that $\mc C=\op{PSh}(S,\mathrm{Ani})$. Then for any $\mc D\in\op{Pr}^L$, we have
	\[\mc C\otimes\mc D\cong\op{PSh}(S,\mc D).\]
	For example, this can be seen directly from the presentation.
\end{example}
The example has supplied us with ``dualizable'' objects.
\begin{definition}[dualizable]
	Let $X$ be an object of a symmetric monoidal category. Then $X$ is \textit{dualizable} if and only if there is an object $X^\lor$ such that $X^\lor\otimes-\cong\op{Mor}(X,-)$.
\end{definition}
\begin{remark}
	It is equivalent to ask that $X\otimes-$ admits a left adjoint of the form $X^\lor\otimes-$. One sees that the definitions are equivalent by moving around adjoints and using the $\otimes$-$\op{Hom}$ adjunction.
\end{remark}
\begin{remark}
	By unwinding the adjunctions, we are provided with a co-evaluation $1\to X\otimes X^\lor$ and an evaluation $X^\lor\otimes X\to1$ for which the composite
	\[X\to X\otimes X^\lor\otimes X\to X\]
	is the identity. This data is basically equivalent to asserting that $X^\lor$ is the dual of $X$.
\end{remark}
\begin{example}
	\Cref{ex:dualizable-object} implies that $\op{PSh}(S;\mathrm{Ani})$ is dualizable for any small $S$.
\end{example}

\subsection{Dualizable Stable Presentable Categories}
As motivation, let's work with some stable categories.
\begin{definition}[stable]
	Define $\op{Pr}^L_{\mathrm{st}}$ to be the category of \textit{stable presentable} categories, meaning that the objects are presentable categories which admit finite limits, finite colimits, a zero object, and all pushout squares are also pullback squares. Morphisms in $\op{Pr}^L_{\mathrm{st}}$ are required to be exact and preserve colimits.
\end{definition}
\begin{remark}
	The internal hom in $\op{Pr}^L_{\mathrm{st}}$ turns out to agree with the internal hom in $\op{Pr}^L$.
\end{remark}
\begin{remark}
	If $S$ is a small stable category, then the free presentable stable $\infty$-category is $\op{Ind}S$; this turns out to be the same as the category of exact presheaves on $S$ (valued in $\mathrm{Sp}$).
\end{remark}
\begin{proposition} \label{prop:compactly-generated-is-dualizable}
	Suppose $\mc C\in\op{Pr}^L_{\mathrm{st}}$ is $\kappa$-compactly generated. Then $\mc C$ is dualizable with dual $\op{Ind}(\mc C^{\kappa,\mathrm{op}})$.
\end{proposition}
\begin{proof}
	Define $\mc C^\lor\coloneqq\op{Ind}(\mc C^{\kappa,\mathrm{op}})$ for brevity. Then one computes that
	\[\mc C^\lor\otimes D=\op{Fun}^{\mathrm{ex}}(\mc C^\kappa,\mc D)\]
	because $\mc C^\lor$ is the free stable category on $\mc C^\kappa$. But now the right-hand side is simply $\op{Fun}^L(\op{Ind}(\mc C^\kappa),\mc D)$, which is $\op{Fun}^L(\mc C,\mc D)$. Thus, we are done!
\end{proof}
Our goal is to classify dualizable presentable categories. For example, it will include the stable ones.
\begin{definition}
	A map $f\colon X\to Y$ of categories is \textit{compact} if and only if the following holds: for all filtered limits $Z=\colim_iZ_i$, there is a dashed arrow making the diagram
	% https://q.uiver.app/#q=WzAsNCxbMCwwLCJcXGRpc3BsYXlzdHlsZVxcY29saW1faVxcb3B7SG9tfShZLFpfaSkiXSxbMCwxLCJcXGRpc3BsYXlzdHlsZVxcY29saW1faVxcb3B7SG9tfShYLFpfaSkiXSxbMSwwLCJcXG9we0hvbX0oWSxaKSJdLFsxLDEsIlxcb3B7SG9tfShYLFopIl0sWzAsMSwiZiIsMl0sWzIsMywiZiJdLFswLDJdLFsxLDNdLFsyLDEsIiIsMSx7InN0eWxlIjp7ImJvZHkiOnsibmFtZSI6ImRhc2hlZCJ9fX1dXQ==&macro_url=https%3A%2F%2Fraw.githubusercontent.com%2FdFoiler%2Fnotes%2Fmaster%2Fnir.tex
	\[\begin{tikzcd}[cramped]
		{\displaystyle\colim_i\op{Hom}(Y,Z_i)} & {\op{Hom}(Y,Z)} \\
		{\displaystyle\colim_i\op{Hom}(X,Z_i)} & {\op{Hom}(X,Z)}
		\arrow[from=1-1, to=1-2]
		\arrow["f"', from=1-1, to=2-1]
		\arrow[dashed, from=1-2, to=2-1]
		\arrow["f", from=1-2, to=2-2]
		\arrow[from=2-1, to=2-2]
	\end{tikzcd}\]
	commute.
\end{definition}
\begin{example}
	One can check that $X$ is compact if and only if $\id_X\colon X\to X$ is compact.
\end{example}
\begin{example}
	If $f$ is a compact morphism, then any composite with $f$ is still compact.
\end{example}
\begin{example}
	A morphism of vector spaces is compact if and only if the image is finite-dimensional.
\end{example}
\begin{lemma} \label{lem:better-compact-morphism}
	Suppose that $\mc C\in\op{Pr}^L_\omega$. Then a map $f\colon X\to Y$ in $\mc C$ is compact if and only if $f$ factors through a compact object.
\end{lemma}
\begin{proof}
	The backward direction follows from the two examples. For the forward direction, suppose $f$ is compact. We may write $Y$ as a filtered colimit $\colim_iZ_i$, where each $Z_i$ is compact. Then compactness of $f$ implies that $f$ factors through $Z_i$ for some $i$.
\end{proof}
\begin{definition}[compactly assembled]
	An object $X$ of a presentable category $\mc C$ is \textit{basic nuclear} if and only if $X$ is a colimit of a diagram of the form
	\[X_0\to X_1\to X_2\to\cdots,\]
	where each morphism $X_i\to X_{i+1}$ is compact. The category $\mc C$ is \textit{compactly assembled} if and only if it is generated (by colimits) by basic nuclear objects.
\end{definition}
\begin{example}
	Let $\mc C$ be the poset category of open subsets of $\RR$. By \Cref{lem:better-compact-morphism}, $U\subseteq V$ is compact if and only if there is a compact set between them. It turns out that $\mc C$ is compactly assembled.
\end{example}
\begin{proposition}
	Basic nuclear objects of a presentable category are $\aleph_1$-compact.
\end{proposition}
\begin{proof}
	Write a basic nuclear object $X$ as promised as $\colim_iX_i$. Now, suppose that $Z=\colim_nZ_n$ is an $\aleph_1$-filtered colimit. We now compute
	\begin{align*}
		\op{Hom}(X,Z) &= \op{Hom}\left(\colim_iX_i,\colim_nZ_n\right) \\
		&= \lim_i\op{Hom}(X_i,\colim_n Z_n) \\
		&= \lim_i\colim_n\op{Hom}(X_i,Z_n).
	\end{align*}
	Now, one can swap the limit and colimit here using the maps
	\[\op{Hom}(X_i,\colim Z_n)\to\colim\op{Hom}(X_i,Z_n)\]
	given by compactness. We conclude that
	\[\op{Hom}(X,Z)=\colim_n\op{Hom}(X,Z_n),\]
	as desired.
\end{proof}
\begin{remark}
	If $\mc C$ is compactly generated, then the converse is also true.
\end{remark}
Here is our main theorem.
\begin{theorem} \label{thm:dualizable-grab-bag}
	Let $\mc C$ be a presentable stable category. Then the following are equivalent.
	\begin{listroman}
		\item $\mc C$ is dualizable.
		\item $\mc C$ is retract of a compactly generated category.
		\item $\mc C$ is $\aleph_1$-presentable, and the natural map $\colim\colon\op{Ind}(\mc C^{\aleph_1})\to\mc C$ admits a left adjoint.
		\item $\mc C$ is compactly assembled.
	\end{listroman}
\end{theorem}
\begin{proof}[Sketch]
	Here we go.
	\begin{itemize}
		\item For (iii) implies (ii), being a fully faithful right adjoint means that the map is localization, so admitting a left adjoint induces a splitting from the localization.
		\item For (ii) implies (i), note that compactly generated categories are dualizable, and dualizability can be checked to be stable under taking retracts.
		\item For (i) implies (iii), again use the fact that $\op{Ind}(\mc C^{\aleph_1})\to\mc C$ is a localization. Thus,
		\[\mc C^\lor\otimes\op{Ind}(\mc C^{\aleph_1})\to\mc C^\lor\otimes\mc C\]
		is also a localization. This is equivalent to
		\[\op{Fun}^L\left(\mc C,\op{Ind}(\mc C^{\aleph_1})\right)\to\op{Fun}^L(\mc C,\mc C)\]
		being a localization. Thus, we may lift morphisms; lifting $\id_{\mc C}$ produces the needed adjunction.
		\item For (iv) implies (iii), we need to define a map $\mc C\to\op{Ind}\left(\mc C^{\aleph_1}\right)$. It is enough to define the map on basic nuclear objects, for which send $X=\colim_iX_i$ to the corresponding colimit in $\op{Ind}\left(\mc C^{\aleph_1}\right)$. The main check is that
		\[\op{Hom}_{\mc C}\left(\colim_iX_i,\colim_nZ_n\right)=\op{Hom}_{\op{Ind}\left(\mc C^{\aleph_1}\right)}\left(\colim_iX_i,\colim_nZ_n\right),\]
		which is purely formal.
		\item For (iii) implies (iv), reverse the argument of the previous point to write any $X\in\mc C^{\aleph_1}$ to some colimit of the form $\colim X_i$ in $\op{Ind}\left(\mc C^{\aleph_1}\right)$. One finds that the maps $X_i\to X$ are compact, which is enough.
		\qedhere
	\end{itemize}
\end{proof}
Now that we understand dualizable objects, we feel comfortable making them a category.
\begin{definition}
	The category $\left(\op{Pr}^L_{\mathrm{st}}\right)^{\mathrm{dual}}$ is the subcategory of dualizable $\mc C\in\op{Pr}^L_{\mathrm{st}}$, where morphisms are functors $\mc C\to\mc D$ whose right adjoints admit right adjoints.
\end{definition}
\begin{remark}
	If $\mc C$ and $\mc D$ are compactly generated, then the presence of all these right adjoints means that compact objects are sent to compact objects.
\end{remark}
\begin{remark} \label{rem:splitting-from-ind}
	It turns out that there is a functorial splitting of $\mc C\to\mc C$ (for dualizable $\mc C$) through $\op{Ind}\left(\mc C^{\aleph_1}\right)$.
\end{remark}
\begin{proposition}
	The forgetful functor to $\left(\op{Pr}^L_{\mathrm{st}}\right)^{\mathrm{dual}}\to\op{Pr}^L_{\mathrm{st}}$ commutes with colimits.
\end{proposition}
\begin{proof}
	Use the functorial splitting of \Cref{rem:splitting-from-ind}. The moral is that we should check that colimits of dualizable categories continue to be dualizable. Well, any such colimit $\colim_n\mc C_n$ admits a factorization
	\[\colim_n\mc C_n\to\colim_n\op{Ind}\mc C_n^{\aleph_1}\to\colim_n\mc C_n,\]
	but the middle term is $\op{Ind}\colim_n\mc C_n^{\aleph_1}$. In particular, the middle term can be seen to be dualizable, and \Cref{thm:dualizable-grab-bag} tells us that the retract $\colim_n\mc C_n$ is still compactly assembled.
\end{proof}
Here are some more facts, whose proofs we omit.
\begin{proposition}
	The category $\left(\op{Pr}^L_{\mathrm{st}}\right)^{\mathrm{dual}}$ is $\aleph_1$-presentable.
\end{proposition}
\begin{proposition}
	Fix an uncountable regular cardinal $\kappa$ and some dualizable stable presentable $\mc C$. Then the following are equivalent.
	\begin{listroman}
		\item $\mc C$ is $\kappa$-compact in $\left(\op{Pr}^L_{\mathrm{st}}\right)^{\mathrm{dual}}$.
		\item $\mc C$ is $\kappa$-compact in $\op{Pr}^L_{\mathrm{st},\kappa}$.
		\item $\mc C$ is dualizable in $\op{Pr}^L_{\mathrm{st},\kappa}$.
	\end{listroman}
\end{proposition}

\subsection{Dualizable \texorpdfstring{$\mc V$}{ V}-Linear Categories}
Throughout, we fix some commutative algebra object $\mc V\in\op{Pr}^L_{\aleph_1}$.
\begin{notation}
	Define $\mathrm{Pr}^L_{\mc V}\coloneqq\mathrm{Mod}_{\mc V}(\mathrm{Pr}^L)$.
\end{notation}
\begin{definition}[internal left adjoint]
	A map $f\colon M\to N$ in $\op{Pr}^L_{\mc V}$ is an \textit{internal left adjoint} if and only if its right adjoint $f^R$ commutes with colimits, and the induced map
	\[v\otimes f^R(x)\to f^R(v\otimes x)\]
	is an equivalence for all $v\in\mc V$ and $x\in N$.
\end{definition}
\begin{definition}[atomic]
	Fix $\mc M\in\op{Pr}^L_{\mc V}$. Then $X\in\mc M$ is \textit{$\mc V$-atomic} if and only if $x\otimes-\colon\mc V\to\mc M$ is an internal left adjoint. Let $\mc M^{\mathrm{at}}$ denote the full subcategory of atomic objects.
\end{definition}
\begin{example}
	By expanding out the definition, we see that if $\mc V=D(\mathbb S)$, then $\mc V$-atomic is equivalent to compact.
\end{example}
\begin{example}
	If $X\in\mc V$, then $X$ is atomic
\end{example}
Thus, we hope that we can replace the word ``compact'' with the word ``atomic'' everywhere in the previous subsection and prove some nice theorems.
\begin{definition}[atomically generated]
	A category $\mc M\in\op{Pr}^L_{\mc V}$ is \textit{atomically generated} if and only if $\mc M$ is generated under colimits (and tensor products with $\mc V$) by atomic objects.
\end{definition}
\begin{definition}
	For any small $\mc V$-enriched category $S$, define $P_{\mc V}(S)$ to be the free presentable $\mc V$-module on $S$.
\end{definition}
\begin{remark}
	If $\mc M$ is $\mc V$-atomically generated, then it turns out that $\mc M\cong P_{\mc V}(\mc M^{\mathrm{at}})$.
\end{remark}
\begin{proposition}
	If $\mc M$ is $\mc V$-atomically generated, then $\mc M$ is dualizable, and its dual is $P_V(\mc M^{\mathrm{at},\mathrm{op}})$.
\end{proposition}
\begin{proof}
	We rewrite the proof of \Cref{prop:compactly-generated-is-dualizable}: indeed,
	\begin{align*}
		\mc M^\lor\otimes\mc N &\cong P_{\mc V}\left(\mc M^{\mathrm{at},\mathrm{op}}\right)\otimes\mc N \\
		&\cong \op{Fun}_{\mc V}\left(\mc M^{\mathrm{at}},\mc N\right) \\
		&\cong \op{Fun}_{\mc V}^L\left(P_V(\mc M^{\mathrm{at}}),\mc N\right) \\
		&\cong \op{Fun}_{\mc V}^L(\mc M,\mc N),
	\end{align*}
	as desired.
\end{proof}
\begin{theorem}
	Fix $\mc M\in\op{Pr}^L_{\mc V}$. Then the following are equivalent.
	\begin{listroman}
		\item $\mc M$ is dualizable.
		\item $\mc M$ is a retract of a $\mc V$-atomically generated category $\mc N\in\op{Pr}^L_{\mc V}$.
		\item $\mc M$ is $\aleph_1$-compactly generated, and the natural map $P_V(\mc M^{\aleph_1})\to\mc M$ admits a left adjoint.
	\end{listroman}
\end{theorem}
\begin{proof}
	Rewrite the proof of \Cref{thm:dualizable-grab-bag}; note that we are avoiding part (iv). I didn't understand the proof the first time, so I will not bother to rewrite it.
\end{proof}
\begin{remark}
	For (iii), the proof shows that the left adjoint is internal.
\end{remark}
\begin{definition}
	Let $\left(\op{Pr}^L_{\mc V}\right)^{\mathrm{dual}}$ be the subcategory of $\op{Pr}^L_{\mc V}$ whose objects are dualizable and whose morphisms are $\mc V$-linear internal left adjoints.
\end{definition}
\begin{remark}
	It turns out that $\mc V$-linear internal left adjoints preserves compact objects.
\end{remark}
\begin{remark}
	One has the same functorial splitting
	\[\mc M\to P_{\mc V}\left(\mc M^{\aleph_1}\right)\to\mc M\]
	in $\op{Pr}^L_{\mc V}$ as before.
\end{remark}
One has the same facts as before. Here is one of them.
\begin{proposition}
	The forgetful functor $\left(\op{Pr}^L_{\mc V}\right)^{\mathrm{dual}}\to\op{Pr}^L_{\mc V}$ commutes with colimits, and $\left(\op{Pr}^L_{\mc V}\right)^{\mathrm{dual}}$ is $\aleph_1$-present\-able.
\end{proposition}

\end{document}